%
%  chapter1.tex — Introdução
%
\chapter{Introdução}
\label{cha:introducao}

\section{Contexto e Motivação}
\label{sec:contexto}

A transição para cidades inteligentes constitui um dos maiores desafios socioeconómicos da atualidade. Este processo exige não apenas infraestruturas tecnológicas avançadas, mas também cidadãos informados, capacitados e participativos. A Aliança Universitária U!REKA, no âmbito da qual o presente projeto se insere, tem como missão acelerar precisamente esta transição, promovendo soluções inovadoras que articulem tecnologia, sustentabilidade e educação~\cite{ureka2024}.

Neste quadro, a literacia cívica e a educação para a sustentabilidade emergem como alavancas fundamentais. A capacidade de os cidadãos adquirirem novas competências de forma contínua e acessível representa um vetor estratégico para a coesão social e para a adaptação às exigências de uma sociedade em permanente transformação. Este imperativo encontra enquadramento direto nos Objetivos de Desenvolvimento Sustentável (ODS) das Nações Unidas, nomeadamente no ODS~4 (\textit{Educação de Qualidade}), ODS~11 (\textit{Cidades e Comunidades Sustentáveis}) e ODS~13 (\textit{Ação Climática}), que apelam à participação ativa dos cidadãos na construção de sociedades mais resilientes e sustentáveis~\cite{un2015sdg}.

A motivação central do \textit{Play4Change} assenta na premissa de que a mudança genuína começa pelo conhecimento. Uma plataforma de aprendizagem acessível, personalizada e gamificada pode capacitar os cidadãos para agirem como agentes de transformação nas suas comunidades, alinhando literacia cívica e ação climática com os objetivos de desenvolvimento sustentável.

\section{Problema}
\label{sec:problema}

Apesar da proliferação de plataformas de aprendizagem online, persistem lacunas significativas no que respeita à educação cívica contextualizada e à promoção de competências direcionadas para a sustentabilidade urbana. As soluções existentes carecem frequentemente de: (i)~mecanismos de personalização efetivos que respondam às dificuldades individuais de cada aprendente; (ii)~modelos de reconhecimento de competências que sejam simultaneamente granulares e verificáveis; e (iii)~integração com os desafios específicos das cidades inteligentes.

Adicionalmente, a motivação intrínseca para a aprendizagem continuada representa um obstáculo recorrente em contextos de educação informal. A ausência de elementos de \textit{engagement} e reconhecimento progressivo contribui para taxas de abandono elevadas nas plataformas convencionais.

\section{Objetivos}
\label{sec:objetivos}

O presente trabalho tem como objetivo principal conceber, implementar e avaliar o \textit{Play4Change}, uma plataforma de aprendizagem adaptativa e gamificada orientada para cidadãos, no âmbito da missão U!REKA.

Os objetivos específicos são:
\begin{enumerate}
  \item Desenvolver um cliente móvel que permita aos cidadãos realizar tarefas de aprendizagem diárias de forma gamificada, adquirindo microcompetências progressivas reconhecidas por \textit{badges};
  \item Implementar um portal de administração web que habilite gestores a criar e gerir tópicos formativos;
  \item Construir um servidor \textit{stateless} que coordene todos os componentes da plataforma de forma escalável;
  \item Integrar um agente de inteligência artificial capaz de gerar conteúdo formativo automaticamente, detetar dificuldades de aprendizagem em tempo real e propor percursos alternativos personalizados;
  \item Implementar um mecanismo de reutilização de percursos adaptativos entre cidadãos com perfis de dificuldade semelhantes, potenciando a eficácia coletiva do sistema.
\end{enumerate}

\section{Contribuições}
\label{sec:contribuicoes}

As principais contribuições deste trabalho são:
\begin{itemize}
  \item Um modelo de microcompetências gamificadas aplicado ao domínio das cidades inteligentes e da sustentabilidade urbana;
  \item Uma arquitetura de geração adaptativa de conteúdo por IA, capaz de detetar dificuldades de aprendizagem em tempo real e propor percursos personalizados;
  \item Um mecanismo de reutilização de percursos por semelhança vetorial --- a contribuição técnica central: os percursos adaptativos são indexados em \textit{pgvector} e recuperados para perfis análogos, acumulando progressivamente conhecimento pedagógico coletivo.
\end{itemize}

\section{Estrutura do Documento}
\label{sec:estrutura}

O presente documento encontra-se organizado da seguinte forma. O Capítulo~\ref{cha:enquadramento} apresenta o enquadramento teórico e o estado da arte, contextualizando o projeto face às soluções e tecnologias existentes. O Capítulo~\ref{cha:solucao} descreve a solução proposta, detalhando a arquitetura e os componentes principais da plataforma. O Capítulo~\ref{cha:implementacao} descreve as decisões de implementação mais relevantes e o \textit{stack} tecnológico adotado. O Capítulo~\ref{cha:avaliacao} apresenta a metodologia de avaliação e os resultados obtidos. Por fim, o Capítulo~\ref{cha:conclusoes} sintetiza as conclusões do trabalho e aponta direções para investigação futura.
